% I livelli che la risoluzione mereologica attraversa: le classi degli oggetti che un
% singolo punto costituisce, dal piu' esteso al piu' specifico.
% Incluso da chapters/introduction.tex (sottosezione Il framework).
\begin{tikzpicture}[
  font=\footnotesize,
  node distance=6mm,
  cls/.style={draw, semithick, rounded corners=3pt, fill=black!6,
              inner sep=3pt, minimum height=5.5mm, minimum width=32mm},
  leaf/.style={cls, double, fill=black!16},
  livello/.style={-{Stealth[length=4pt]}, semithick, black!55},
]
  \node[cls]              (build)  {\texttt{base:Building}};
  \node[cls,  below=of build]  (fac) {\texttt{base:Facade}};
  \node[cls,  below=of fac]    (col) {\texttt{base:Column}};
  \node[leaf, below=of col]    (cap) {\texttt{base:Capital}};

  \draw[livello] (build) -- (fac);
  \draw[livello] (fac)   -- (col);
  \draw[livello] (col)   -- (cap);

  % l'asse, non gli archi: non sono triple, sono livelli di risoluzione
  \draw[-{Stealth[length=5pt]}, semithick]
    ($ (build.north east) + (5mm,0) $) -- ($ (cap.south east) + (5mm,0) $);
  \node[right=1mm, font=\scriptsize, align=left, anchor=north west]
    at ($ (build.north east) + (5mm,0) $) {risoluzione\\mereologica};
\end{tikzpicture}
